% from ega0-14-fr.tex
\oldpage[0\textsubscript{IV-1}]{10}
\begin{proposition}[14.2.5]
\label{0.14.2.5-fr}
Si $(Y_i)_{1\leq i\leq n}$ est une famille finie de parties fermées d'un
espace topologique $X$, et $Y$ la réunion de cette famille, on a
\[
  \operatorname{codim}(Y,X)=
  \inf_i\bigl(\operatorname{codim}(Y_i,X)\bigr).
  \tag{14.2.5.1}\label{0.14.2.5.1-fr}
\]
\end{proposition}

\begin{proof}
En effet, toute composante irréductible de l'un des $Y_i$ est contenue dans
une composante irréductible de $Y$, et inversement toute composante
irréductible de $Y$ est aussi composante irréductible de l'un des $Y_i$
($0_{\mathrm I}$, 2.1.1)~; la conclusion résulte donc de la définition
(14.2.1) et de l'inégalité (14.2.2.3).
\end{proof}

\begin{corollary}[14.2.6]
\label{0.14.2.6-fr}
Soient $X$ un espace topologique, $Y$ un sous-espace fermé localement
noethérien de $X$.
\begin{enumerate}
  \item[\rm(i)] Pour tout $x\in X$, il n'existe qu'un nombre fini de
  composantes irréductibles $Y_i$ ($1\leq i\leq n$) de $Y$ contenant $x$,
  et l'on a
  $\operatorname{codim}_x(Y,X)=
  \inf_i(\operatorname{codim}(Y_i,X))$.
  \item[\rm(ii)] La fonction
  $x\mapsto\operatorname{codim}_x(Y,X)$ est semi-continue inférieurement
  dans $X$.
\end{enumerate}
\end{corollary}

\begin{proof}
En effet, par hypothèse il y a un voisinage ouvert $U_0$ de $x$ dans $X$ tel
que $Y\cap U_0$ soit noethérien, donc n'ait qu'un nombre fini de composantes
irréductibles, qui sont traces sur $U_0$ de composantes irréductibles de $Y$~;
\emph{a fortiori} il n'y a qu'un nombre fini de composantes irréductibles
$Y_i$ ($1\leq i\leq n$) de $Y$ contenant $x$, et on peut, en remplaçant $U_0$
par un voisinage ouvert $U\subset U_0$ de $x$ ne rencontrant aucun des $Y_j$
qui ne contiennent pas $x$, supposer que les $Y_i\cap U$ sont les composantes
irréductibles de $Y\cap U$~; pour tout voisinage ouvert $V\subset U$ de $x$
dans $X$, les $Y_i\cap V$ sont alors les composantes irréductibles de
$Y\cap V$, et (14.2.3) montre alors que
$\operatorname{codim}(Y_i,X)=\operatorname{codim}(Y_i\cap V,V)$, ce qui
prouve (i). En outre, pour tout $x'\in U$, les composantes irréductibles de
$Y$ contenant $x'$ sont certaines des $Y_i$, donc
$\operatorname{codim}_{x'}(Y,X)\geq\operatorname{codim}_x(Y,X)$, ce qui
démontre (ii).
\end{proof}

\subsection{La condition des chaînes}
\label{subsection:0.14.3-fr}

\begin{env}[14.3.1]
\label{0.14.3.1-fr}
Dans un espace topologique $X$, nous dirons qu'une chaîne
$Z_0\subset Z_1\subset\cdots\subset Z_n$ de parties fermées irréductibles est
\emph{saturée} s'il n'existe pas de partie fermée irréductible $Z'$ distincte
des $Z_i$ et telle que $Z_k\subset Z'\subset Z_{k+1}$ pour un indice $k$.
\end{env}

\begin{proposition}[14.3.2]
\label{0.14.3.2-fr}
Soit $X$ un espace topologique tel que, pour deux parties fermées
irréductibles quelconques $Y$, $Z$ de $X$ telles que $Y\subset Z$, on ait
$\operatorname{codim}(Y,Z)<+\infty$. Les deux conditions suivantes sont
équivalentes~:
\begin{enumerate}
  \item[\rm(a)] Deux chaînes saturées de parties fermées irréductibles de $X$,
  ayant mêmes extrémités, ont même longueur.
  \item[\rm(b)] Si $Y$, $Z$, $T$ sont trois parties fermées irréductibles de
  $X$ telles que $Y\subset Z\subset T$, on a
  \[
    \operatorname{codim}(Y,T)=
    \operatorname{codim}(Y,Z)+\operatorname{codim}(Z,T).
    \tag{14.3.2.1}\label{0.14.3.2.1-fr}
  \]
\end{enumerate}
\end{proposition}

\begin{proof}
Il est immédiat que a) entraîne b). Réciproquement, supposons b) vérifiée, et
démontrons que si deux chaînes saturées de mêmes extrémités ont pour longueurs
$m$ et $n\leq m$, on a nécessairement $m=n$. Raisonnons par récurrence sur
$n$, la proposition étant évidente pour $n=1$. Supposons donc $n>1$, $n<m$,
et soit $Z_0\subset Z_1\subset\cdots\subset Z_n$ une
